% ==================================================
% 2^n and Its Role in Discrete Calculus
% Discrete Calculus — Special Functions
% ==================================================

\section{The Function \(2^n\) in Discrete Calculus}
\label{subsec:two-to-the-n}

Among all discrete exponentials, the case \(a = 1\) — giving \((1+1)^n = 2^n\)
— is perhaps the most important. It appears in combinatorics as the count
of subsets, in probability as the sample space of fair coin sequences, in
computer science as the natural measure of binary information, and in discrete
calculus as a canonical test function. This section develops its properties
systematically.

\section*{Differencing \(2^n\)}

Setting \(a = 1\) in the discrete exponential rule
(Proposition~\ref{prop:discrete-exponential-rule}) gives immediately:

\begin{tcolorbox}[colback=propbox, colframe=propborder, arc=2pt,
  left=6pt, right=6pt, top=4pt, bottom=4pt,
  title={\small\textbf{Proposition (\(\Delta 2^n = 2^n\))}},
  fonttitle=\small\bfseries]
\[
\Delta\, 2^n = 2^n.
\]
\end{tcolorbox}
\label{prop:delta-two-n}

\begin{remark*}[Fully quantified form]
\(\forall n \in \mathbb{Z}:\; 2^{n+1} - 2^n = 2^n\).
\end{remark*}

\begin{proof}
\[
\Delta\, 2^n = 2^{n+1} - 2^n = 2\cdot 2^n - 2^n = (2-1)\cdot 2^n = 2^n. \qedhere
\]
\end{proof}

\begin{remark*}[Self-similarity]
This says that \(2^n\) is its own forward difference. In continuous calculus,
the analogous property is \(\tfrac{d}{dx}e^x = e^x\): the exponential is its
own derivative. The function \(2^n\) is the discrete analogue of \(e^x\)
in the sense that \(\Delta(2^n) = 2^n\). More precisely, \(2^n = e^{n \ln 2}\),
and the relationship \(\Delta(2^n) = 2^n\) corresponds to the special case
\(a = 1\) of \(\Delta(1+a)^n = a(1+a)^n\).
\end{remark*}

\section*{Antiderivative and Geometric Sum}

Since \(\Delta(2^n) = 2^n\), a discrete antiderivative of \(2^n\) is \(2^n\)
itself (up to a constant).

\begin{tcolorbox}[colback=propbox, colframe=propborder, arc=2pt,
  left=6pt, right=6pt, top=4pt, bottom=4pt,
  title={\small\textbf{Proposition (Antiderivative of \(2^n\))}},
  fonttitle=\small\bfseries]
A discrete antiderivative of \(2^n\) is \(2^n\) itself:
\[
\Delta(2^n) = 2^n
\qquad \text{so} \qquad
\sum_{k=0}^{n-1} 2^k = 2^n - 1.
\]
\end{tcolorbox}
\label{prop:antideriv-two-n}

\begin{remark*}[Fully quantified form]
\(\forall n \in \mathbb{N}:\;
\displaystyle\sum_{k=0}^{n-1} 2^k = 2^n - 1\).
\end{remark*}

\begin{proof}
Applying the Fundamental Theorem of Finite Calculus
(Theorem~\ref{thm:ftfc}) with \(F(k) = 2^k\) and \(f(k) = \Delta F(k) = 2^k\):
\[
\sum_{k=0}^{n-1} 2^k = F(n) - F(0) = 2^n - 2^0 = 2^n - 1. \qedhere
\]
\end{proof}

\begin{remark*}[Verification by expansion]
Expanding: \(1 + 2 + 4 + 8 + \cdots + 2^{n-1}\). Multiply the partial sum
\(S = \sum_{k=0}^{n-1}2^k\) by \(2\) to get \(2S = \sum_{k=1}^{n}2^k\).
Subtracting: \(2S - S = 2^n - 1\), so \(S = 2^n - 1\). This is the
classical derivation of the geometric series formula; the discrete calculus
approach via antidifferentiation gives the same result mechanically.
\end{remark*}

\section*{Higher Differences of \(2^n\)}

\begin{proposition}[Higher differences of \(2^n\)]
\label{prop:higher-diff-two-n}
For all \(k \ge 0\),
\[
\Delta^k\, 2^n = 2^n.
\]
\end{proposition}

\begin{remark*}[Fully quantified form]
\(\forall k \in \mathbb{N},\; \forall n \in \mathbb{Z}:\; \Delta^k(2^n) = 2^n\).
\end{remark*}

\begin{proof}
By induction. The base case \(k = 0\) is \(\Delta^0(2^n) = 2^n\) by
definition. For the inductive step, assume \(\Delta^k(2^n) = 2^n\). Then
\[
\Delta^{k+1}(2^n) = \Delta(\Delta^k(2^n)) = \Delta(2^n) = 2^n. \qedhere
\]
\end{proof}

\begin{remark*}[Consequence for difference tables]
Since every finite difference of \(2^n\) is \(2^n\) again, the difference
table of \(2^n\) consists entirely of the sequence \(2^n\) at every level.
This is the discrete analogue of the fact that the Taylor series of \(e^x\)
has the same coefficient pattern at every derivative level. In particular,
\(2^n\) is \emph{not} a polynomial — a polynomial of degree \(d\) would
reach \(0\) at the \((d+1)\)-th difference row. The non-termination of
the difference table is the signature of a genuinely transcendental discrete
function.
\end{remark*}

\section*{\(2^n\) in Combinatorics and the Newton Expansion}

The Binomial Theorem (Theorem~\ref{thm:binomial}) gives a direct connection:
setting \(x = y = 1\) yields
\[
\sum_{k=0}^{n} \binom{n}{k} = 2^n.
\]
This counts the total number of subsets of an \(n\)-element set: there are
\(\binom{n}{k}\) subsets of size \(k\), and the sum over all \(k\) gives
\(2^n\).

More deeply, the Newton expansion (Theorem~\ref{thm:newton-expansion}) of \(2^n\)
in the falling factorial basis is:
\[
2^n = \sum_{k=0}^{n} \binom{n}{k} = \sum_{k=0}^{n} \frac{n^{\underline{k}}}{k!},
\]
since \(\Delta^k(2^n)\big|_{n=0} = 2^0 = 1\) for every \(k\). This expresses
\(2^n\) as the sum of all normalised falling factorials up to degree \(n\),
and reveals why binomial coefficients count subsets: they are the coefficients
of \(2^n\) in its own Newton expansion.

\begin{remark*}[Analogy with \(e^x\)]
In continuous calculus, \(e^x = \sum_{k=0}^{\infty} x^k/k!\). The Newton
expansion gives \(2^n = \sum_{k=0}^{n} n^{\underline{k}}/k!\). Both express
a canonical exponential as a sum of normalised basis functions, with all
coefficients equal to \(1\). The analogy is exact: ordinary powers replace
falling factorials, and the infinite series terminates at degree \(n\) in
the discrete case because \(2^n\) has degree exactly \(n\) in the falling
factorial basis.
\end{remark*}

\section*{Discrete Differential Equation}

In continuous calculus, the equation \(f' = f\) (with \(f(0) = 1\)) has the
unique solution \(f(x) = e^x\). In discrete calculus, the analogous
\emph{difference equation} is:
\[
\Delta F(n) = F(n), \qquad F(0) = 1.
\]
That is, \(F(n+1) - F(n) = F(n)\), so \(F(n+1) = 2F(n)\). The unique solution
to this initial value problem is \(F(n) = 2^n\).

\begin{proposition}[Discrete IVP]
\label{prop:discrete-ivp}
The unique function \(F : \mathbb{N} \to \mathbb{R}\) satisfying
\[
F(n+1) = 2\,F(n), \qquad F(0) = 1
\]
is \(F(n) = 2^n\).
\end{proposition}

\begin{remark*}[Fully quantified form]
\(\forall n \in \mathbb{N}:\;
(F(n+1) = 2F(n)\text{ and }F(0)=1) \Rightarrow F(n) = 2^n\).
\end{remark*}

\begin{proof}
Uniqueness: if \(G\) is another solution, then \(G(0) = 1 = F(0)\), and
inductively \(G(n+1) = 2G(n) = 2F(n) = F(n+1)\), so \(G = F\).
Existence: verify directly that \(2^{n+1} = 2\cdot 2^n\) and \(2^0 = 1\).
\end{proof}

\begin{remark*}[Significance]
This characterises \(2^n\) as the unique fixed point of ``multiply by 2''
starting from \(1\). The analogous characterisation of \(e^x\) is: unique
solution to \(f' = f\), \(f(0) = 1\). Both are initial-value problems in
their respective calculi, and both have closed-form solutions that serve as
the canonical exponential of that calculus.
\end{remark*}
